\section{Related Work}
\label{sec:related}

\paragraph{LLM-generated simulators.} Several systems build executable models from text. \otf compiles process-record documents into reusable readers and simulators and accepts a program if it reproduces hidden recorded outcomes~\cite{otftools2026}. \devs generates discrete-event models from specifications and validates their event traces against temporal, causal and semantic constraints derived from the specification~\cite{chen2026devs}. G-Sim proposes simulator structure with an LLM and calibrates parameters against data~\cite{holt2025gsim}. Program-aided reasoning more generally lets a model write code and trusts its execution~\cite{gao2023pal}. Closest in spirit, a validation-aware framework for LLM-driven simulation separates whether a simulation runs from whether its numbers are valid~\cite{yoo2026validation}, and work on LLM agents as suppliers of Modelica models asks what an acceptance test actually certifies~\cite{frochte2026modelica}. These systems judge a model, a run or a trace. We ask a narrower question that arises once a model is reused: whether one answer, for one requested change, can be trusted, and we keep evidence per kind of change across reuse.

\paragraph{Structured requests and abstention.} Turning a natural-language what-if request into a checkable specification has been used to build interactive what-if interfaces, where many generated specifications are wrong even when they compile~\cite{gathani2026praxa}. For database answers consumed as fact, a deterministic kernel that answers only a bounded set of question shapes, and abstains otherwise, keeps a generative front end from choosing the returned value~\cite{wu2026never}. Our contract check plays a similar role for changes, but the decision to trust rests on tests of the program. Selective prediction studies the trade-off between coverage and risk~\cite{elyaniv2010selective,geifman2017selective}, and conformal risk control has been extended to agent trajectories~\cite{opoku2026toolchain}. We borrow the coverage-matched comparison, but our gate is binary, makes no calibration assumption, and claims no guarantee.

\paragraph{Models as judges.} LLMs are widely used to grade other models' outputs~\cite{zheng2023judging}, and they can favor outputs from their own family~\cite{panickssery2024self}, which matters here because the same model families build and judge programs. ReAct interleaves reasoning with tool calls~\cite{yao2023react}; our ReAct judge gives it a code tool and the program. Choosing between text reasoning and program execution per question has also been studied as model selection~\cite{zhao2023selection}. Agreement between independent answerers is a common reliability signal, but errors they share stay invisible to it~\cite{zerhoudi2026wrong}; we include disagreement with a reading answer as a judge for that reason.

\paragraph{Testing simulation programs.} Comparing a model's output with recorded history is a standard validation technique for simulation models~\cite{sargent2013vv}. Our change-specific tests are metamorphic relations~\cite{chen2018metamorphic,segura2016survey}: properties that must hold between runs on related inputs when no oracle for the output exists. One of our IT checks, which removes the dominating machine so that a resource-request test can tell CPU from memory, follows constraint-based test-data generation~\cite{demillo1991constraint}. Fixing exogenous draws so that a factual and a changed run can be compared exactly is the common-random-numbers idea from simulation~\cite{law2015simulation}, and matches the structural view of counterfactuals in which exogenous variables are held fixed~\cite{pearl2009causality}. Process mining and business-process simulation supply the workflow models we use~\cite{vanderaalst2016process,lopezpintado2023prosimos,chapela2024extraneous}.
